\section{Method}
\subsection{Overview}
For a 3D mesh $\mathcal{M} = (V,F)$, \ourmodel decomposes the mesh faces into a small collection of disjoint charts:
\vspace{-0.5em}
\begin{equation}
F = \bigcup_{k=1}^{K} C_k, \quad \text{with} \quad C_i \cap C_j = \varnothing \quad (i \neq j),
\vspace{-1ex}
\end{equation}
where each \textbf{\emph{chart}} $C_k$ is a connected subset of faces. We utilize an Angle-Based Flattening algorithm(ABF++), to flatten each chart onto a 2D plane, yielding mappings:
\vspace{-0.5em}
\begin{equation}
\phi_k : C_k \longrightarrow \mathbb{R}^2,
\vspace{-1ex}
\end{equation}
so that each vertex $v_i \in V$ receives a corresponding 2D UV coordinate $\mathbf{u}_i = \phi_k(v_i) \in \mathbb{R}^2$, determined by the chart it belongs.

The primary challenge lies in decomposing the mesh into charts $\{C_k\}$, a task essentially equivalent to identifying optimal seams for mesh segmentation. Many state-of-the-art methods, such as xatlas and Blender's Smart UV, rely on geometric heuristics—such as bottom-up greedy chart expansion based on face normals or distortion metrics—to generate these charts. However, these strategies typically produce overly fragmented charts and unintuitive or suboptimal boundaries, often splitting semantically coherent regions across multiple charts.

In contrast, \ourmodel adopts a coarse-to-fine, two-stage strategy for mesh decomposition. At a high level, it employs a recent learning-based method called PartField (introduced in Section~\ref{sec:partfield}) to partition mesh faces into semantically coherent \textbf{\emph{parts}}: $F = \bigcup_{m=1}^{M} P_m$, where each part $P_k$ exhibits relatively simple geometry—for instance, cylindrical limbs, spherical toes. Subsequently, we introduce two geometry-based heuristics (detailed in Section~\ref{sec:heuristic}) to further segment each part into a small set of \textbf{\emph{charts}}:
$P_m = \bigcup_{n=1}^{N_m} C_{m,n}$, ensuring each resulting chart $C_{m,n}$ can be flattened onto a 2D plane with minimal distortion.

To accomplish this, \ourmodel employs a top-down recursive decomposition search (detailed in Section~\ref{sec:search}) that minimizes the total number of generated charts while ensuring the distortion for each chart remains below a user-specified threshold $\tau$. Since exhaustive decomposition searches can incur substantial computational overhead, we introduce acceleration and parallelization techniques for efficiency. To ensure a complete and robust pipeline, Section~\ref{sec:runtime} also describes additional preprocessing and postprocessing procedures, including handling non-manifold and multi-component meshes, as well as performing UV packing.

\subsection{Preliminary: PartField}
\label{sec:partfield}

PartField~\cite{liu2025partfield} trains a feed‑forward neural network that takes a 3D mesh as input and predicts a continuous, part-based feature field encoded as a triplane. By leveraging extensive contrastive learning on part-labeled 3D data and large‑scale unlabeled 3D data with 2D pseudo part labels, PartField learns general hierarchical concepts of semantic and geometric parts. For any 3D point, we can obtain its high‑dimensional part feature by interpolating the triplane representation. Points whose features are similar—as measured by cosine similarity—are therefore more likely to belong to the same part.
\subsection{Top-Down Recursive Tree Search}
\label{sec:search}

After obtaining the hierarchical part-based feature field using PartField, we first compute a representative part feature for each triangular face by uniformly sampling $s$ 3D points within the face and averaging their corresponding point features. Next, we construct a hierarchical part tree $\mathcal{T}$ using agglomerative clustering~\cite{johnson1967hierarchical} on these face features. In this hierarchical tree structure, the leaf nodes represent individual triangular faces, and the root node encompasses the entire mesh.

% ------------------------------------------------------------------
\begin{algorithm}[t]
\caption{\textsc{PartTreeSearch}}
\label{alg:search}
\begin{algorithmic}[1]
\Require $\mathcal{P}$ (a PartField subtree node), distortion threshold $\tau$, chart budget $B$
\Ensure A \emph{UVChartSet} $\mathscr{C}=\{C_1, C_2, \cdots, C_k\}$ for $\mathcal{P}$'s mesh, whose charts $C_i$ have distortion $\le \tau$ and whose total chart count $k$ is $\le B$; return $\bot$ if no such set exists.

\vspace{2pt}

\Procedure{PartTreeSearch}{$\mathcal{P}$,\,$\tau$,\,$B$}
    \If{$B < 1$}                          \Comment{budget exhausted}
        \State \Return $\bot$
    \EndIf
% ------------------------------------------------------------------
    \State $\mathcal{H}_1 \gets \Call{GenCandidatesH1}{\mathcal{P}.mesh,10}$ \Comment{Heuristic1} \label{line:heuristic1}
    \ForAll{$\mathscr{C} \in \mathcal{H}_1$} \Comment{$\mathscr{C}$ is one of candidate $UVChartSet$}
        \State \Call{ParameterizeABF}{$\mathscr{C}$}   \Comment{flatten and compute distortion}
    \EndFor
% ------------------------------------------------------------------
    \If{$\min_{\mathscr{C}\in \mathcal{H}_1} \mathscr{C}.\text{dist} > \tau$}            \Comment{no admissible H1 candidate}
        \State $L \gets \Call{PartTreeSearch}{\mathcal{P}.left\_child ,\tau, \infty\,}$
        \State $R \gets \Call{PartTreeSearch}{\mathcal{P}.right\_child ,\tau, \infty\,}$
        \State \Return $L \oplus R$ \Comment{merge chart sets from the subtrees}
    \Else
% ------------------------------------------------------------------
    \State $\mathcal{H}_2 \gets \{\Call{GenCandidateH2}{\mathcal{P}.mesh,\tau}\}$ \Comment{Heuristic2} \label{line:heuristic2}
    \State \Call{ParameterizeABF}{$\mathcal{H}_2[0]$} 
    \State $\mathcal{S} \gets \bigl\{\,\mathscr{C}\in \mathcal{H}_1 \cup \mathcal{H}_2 \mid
            \mathscr{C}.\text{dist}\le\tau \;\land\; \Call{NoOverlap}{\mathscr{C}}\bigr\}$
    \State $\mathscr{C}_{\text{best}} \gets \arg\!\min_{\mathscr{C} \in \mathcal{S}} \mathscr{C}.\text{count}$
% ------------------------------------------------------------------
    \State \Comment{Recurse to determine whether a better solution exists}     \label{line:recursion}
    \State $B' \gets \min(B,\,\mathscr{C}_{\text{best}}.\text{count}-1)$ 
    \State $L \gets \Call{PartTreeSearch}{\mathcal{P}.left ,\tau,B'-1}$
    \State $R \gets \Call{PartTreeSearch}{\mathcal{P}.right,\tau,B'-L.\text{count}}$
    \State $\mathscr{C}_{\text{comb}} \gets L \oplus R$
    \State $is\_valid \gets \mathscr{C}_{\text{comb}}.\text{dist}\le\tau \;\land\; \Call{NoOverlap}{\mathscr{C}_{\text{comb}}}$
    \If{$\mathscr{C}_{\text{comb}}.\text{count}<\mathscr{C}_{\text{best}}.\text{count} \;\land\; is\_valid$}
        \State $\mathscr{C}_{\text{best}}\gets \mathscr{C}_{\text{comb}}$
    \EndIf
    \State \Return $\mathscr{C}_{\text{best}}$
    \EndIf
% ------------------------------------------------------------------

\EndProcedure
\end{algorithmic}

\end{algorithm}
% \vspace{-2ex}


A straightforward approach to mesh decomposition might involve using PartField to directly generate a fixed number of parts and attempting to flatten each to 2D individually. Alternatively, one could adaptively traverse $\mathcal{T}$ from the root downward, checking whether each node’s corresponding geometry can be flattened into a 2D chart with minimal distortion. However, we observe that relying solely on parts generated by PartField often leads to suboptimal results. This limitation arises because PartField primarily focuses on semantic or coarse geometric partitioning and is less effective for finer-scale, UV-related decompositions—such as accurately segmenting cylindrical or spherical regions into charts that minimize distortion. To address this challenge, \ourmodel leverages PartField primarily for high-level semantic decomposition into geometrically simpler subparts and employs two geometry-based heuristics (elaborated in Section~\ref{sec:heuristic}) to further divide these semantically meaningful parts into smaller charts amenable to low-distortion flattening. The search algorithm interleaves these two strategies.

Formally, \ourmodel employs a top-down recursive decomposition search—detailed in Algorithm~\ref{alg:search}—to optimally balance the chart count against distortion constraints. Given a subtree node $\mathcal{P}$ (initially the root node of the tree $\mathcal{T}$), a distortion threshold $\tau$, and a chart budget $B$ (initially set to $\infty$), the algorithm searches for a valid chart decomposition of $\mathcal{P}$'s mesh that satisfies three key conditions: each chart has a distortion of at most $\tau$, no charts overlap in 2D, and the total number of charts does not exceed $B$. If no feasible decomposition exists, the algorithm returns failure ($\bot$).

Specifically, the algorithm begins by generating candidate chart decompositions for node $\mathcal{P}$'s mesh using a primary geometric heuristic (Heuristic1). Each candidate chart set $\mathscr{C} = \{C_i\}$, consisting of up to $t$ charts, is flattened using Angle-Based Flattening (ABF), and its distortion is evaluated. If none of the candidate decompositions with up to $t$ charts from Heuristic1 satisfy the distortion constraint (i.e., $\delta_{\min} > \tau$), the algorithm utilizes the PartField tree to divide the mesh and recursively searches both the left and right child subtrees without budget constraints, then merges their respective optimal chart sets.

Conversely, if an admissible candidate decomposition is found through Heuristic1, we further refine it using a secondary heuristic (Heuristic2) designed to potentially yield fewer charts. Among these candidates, we select the best solution $\mathscr{C}_{\text{best}}$ with the minimal chart count that satisfies the distortion and overlap constraints.

Before accepting this solution, the algorithm performs a final check by recursively exploring the left and right child nodes of the PartField subtree, using a reduced chart budget ($B'$), which is derived from the best candidate’s chart count minus one. If the combined solution from this subtree search yields a valid and superior decomposition (i.e., fewer charts than $\mathscr{C}_{\text{best}}$), it replaces the previously identified best solution. The chart budget $B$ prevents the search from going too deep. As the search progresses, $B$ tightens, and recursion stops when candidates exceed it. In practice, the search rarely goes very deep.

By strategically interleaving hierarchical semantic guidance from PartField with fine-grained geometric heuristics and systematic recursive exploration, our proposed decomposition search achieves semantically coherent, distortion-bounded, and notably compact mesh parameterizations—characterized by a small number of charts.

\subsection{Geometry-Based Part Decomposition}
\label{sec:heuristic}
In this section, we introduce two heuristics, \emph{Normal} and \emph{Merge}, to further decompose a part $P_m$ generated by PartField and exhibiting simple geometry. The goal is to divide $P_m$ into multiple charts, denoted as $P_m = \bigcup_{n=1}^{N_m} C_{m,n}$, such that each chart can be flattened to 2D with low distortion.

The first heuristic, referred to as \emph{Normal} (line~\ref{line:heuristic1} in Algorithm~\ref{alg:search}), is based on face normals. We apply an agglomerative clustering algorithm~\cite{johnson1967hierarchical} to the face normals of $P_m$'s mesh, partitioning its triangle faces into charts, where each chart is composed of connected faces with similar normals. This clustering is performed once for $P_m$, producing $t$ candidate decompositions with 1 to $t$ charts (we use $t=10$ in our experiments). For each candidate decomposition, we apply the Angle-Based Flattening (ABF) algorithm to flatten the charts and evaluate distortion. Since ABF aims to preserve angles, we quantify distortion using an area stretch metric, defined as:
\vspace{-0.3em}
\begin{equation}
\label{equ:distortion}
\small
\text{distortion}(\mathscr{C}) = \max_{C \in \mathscr{C}} \left( \frac{1}{|C|} \sum_{f \in C} \max\left( \text{stretch}(f), \frac{1}{\text{stretch}(f)} \right) \right),
\end{equation}
\vspace{-0.3em}
where $C$ denotes a single chart composed of multiple connected faces, and $\mathscr{C}$ denotes the set of all charts in the decomposition. The stretch of a triangle face $f$ is computed as:
\vspace{-0.4em}
% \begin{equation}
% \small
% \text{stretch}(f) = \frac{\text{area}{2D}(f)}{\text{area}{3D}(f)} \Bigg/ \left( \frac{1}{|C|} \sum_{f' \in C} \frac{\text{area}{2D}(f')}{\text{area}{3D}(f')} \right).
% \end{equation}
\begin{equation}
\small
\text{stretch}(f) = \frac{\text{area}{2D}(f)}{\text{area}{3D}(f)} \Bigg/ \left( \frac{\sum_{f' \in C} \text{area}{2D}(f')}{\sum_{f' \in C} \text{area}{3D}(f')} \right).
\end{equation}
Note that both PartField and the \emph{Normal} heuristic employ the agglomerative clustering algorithm to group faces into parts or charts. The key difference lies in the features used: PartField utilizes learned high-level part features, while heuristic \emph{Normal} relies on low-level geometric face normals. These two strategies are thus consistent in spirit and complementary in practice.

The \emph{Normal} heuristic is simple, fast, and often yields satisfactory results. However, we also propose a second, more computationally expensive heuristic, called \emph{Merge} (line~\ref{line:heuristic2} in Algorithm~\ref{alg:search}), which may produce decompositions with fewer charts. Given a part $P_m$, the \emph{Merge} heuristic begins by computing its oriented bounding box (OBB). It then assigns each triangle face a label from 1 to 6, corresponding to the OBB face normal with which the triangle’s normal is most closely aligned. Using both face connectivity and these labels, we segment the faces into multiple connected components. These components are then sorted by size (small to large), and we iteratively attempt to merge each component with its adjacent neighbors, starting from the one with the longest shared edges. For each merge attempt, we temporarily merge two components and apply ABF to flatten the combined chart. The merge is accepted if the resulting chart satisfies the distortion threshold and is free of overlaps. This merging process continues until no further valid merges can be made, after which the final chart set is returned. Since the \emph{Merge} heuristic often begins with many small components and performs multiple ABF calls during its iterative merging process, it is significantly more expensive than the \emph{Normal} heuristic. However, it may yield better decompositions with fewer charts. Therefore, we only invoke the \emph{Merge} heuristic when the \emph{Normal} heuristic returns a valid (admissible) decomposition.

\subsection{Runtime Optimization and Pre- and Post-processing}
\label{sec:runtime}

\textbf{Runtime Optimization} To ensure efficiency despite the computational cost of recursive decomposition and repeated ABF invocations, we adopt two key strategies during the decomposition process. First, we parallelize all recursive calls to the left and right subtrees during the top-down search, allowing the algorithm to exploit multi-core processing and significantly accelerate the overall decomposition. Second, to avoid the high cost of repeatedly invoking Angle-Based Flattening (ABF) on dense meshes, we employ a GPU-accelerated mesh simplification algorithm~\cite{oh2025pamo} to generate low-resolution approximations of candidate charts. During simplification, the chart boundary is kept fixed to preserve the geometric structure relevant to UV mapping. These simplified charts are used to estimate distortion metrics quickly (i.e., surrogate distortion) during intermediate evaluations. Once the final chart set is determined, ABF is applied to the original high-resolution mesh to produce accurate UV coordinates and distortion measurements.

\noindent \textbf{Non-Manifold and Multi-Component Meshes} Since the ABF algorithm assumes each input chart is a manifold and connected surface, additional processing is required when handling non-manifold or multi-component meshes. For non-manifold meshes, we detect all non-manifold edges—i.e., edges shared by more than two faces—and resolve them by duplicating the shared vertices and splitting each such edge into $N-1$ distinct edges, where $N$ is the number of incident faces, thereby converting the structure into a manifold form suitable for flattening. For meshes containing multiple connected components, we initially proceed with the proposed PartField-guided hierarchical decomposition. If a part mesh consists of disconnected components, we skip heuristic-based decomposition at that level and instead recursively explore the left and right subtrees of the PartField hierarchy. However, if PartField fails to further segment the multiple components after reaching a predefined recursion depth, we fall back to applying the geometric heuristics and ABF flattening to each connected component at that level individually, in order to avoid excessively fragmented decomposition.

\noindent \textbf{UV Packing} While our primary focus is on decomposing 3D meshes into charts and generating corresponding low-distortion 2D parameterizations, our method is fully compatible with a variety of existing UV packing algorithms. A distinguishing feature of our approach is that chart decompositions are grouped based on semantically meaningful parts, enabling more structured and application-aware packing strategies. 
For example, charts belonging to the same part can be grouped together during packing. 
Alternatively, part groups can be packed into an arbitrary number of UV atlas squares (e.g., multiple $[0,1]^2$ spaces) with a semantically balanced distribution across atlases.
This semantic hierarchy not only improves organizational clarity but also benefits downstream applications such as texture painting or editing, where charts belonging to the same part remain spatially close and are easier to manipulate collectively.